% Generated draft: requires mathematical and source review.
\documentclass[11pt]{article}
\usepackage[margin=25mm]{geometry}
\usepackage{fontspec}
\setmainfont{Latin Modern Roman}
\usepackage{amsmath,amssymb,mathtools,unicode-math}
\setmathfont{Latin Modern Math}
\usepackage{enumitem,needspace,xurl,hyperref}
\setlength{\parindent}{0pt}
\setlength{\parskip}{5pt}
\setlength{\emergencystretch}{3em}
\newcommand{\N}{\mathbb{N}}
\newcommand{\Z}{\mathbb{Z}}
\newcommand{\Q}{\mathbb{Q}}
\newcommand{\R}{\mathbb{R}}
\newcommand{\C}{\mathbb{C}}
\newcommand{\F}{\mathbb{F}}
\newcommand{\A}{\mathbb{A}}
\newcommand{\PP}{\mathbb{P}}
\newcommand{\E}{\mathbb{E}}
\newcommand{\eps}{\varepsilon}
\newcommand{\dd}{\,\mathrm{d}}
\begin{document}
\section*{UnsolvedMath ID: 30002677}
\textbf{Source problem code:} catalogue entry
\subsection{Definitions and mathematical statement}
A stratified space is a space \(X\) decomposed into smooth strata \(S_\alpha\), written \(X=\bigsqcup_\alpha S_\alpha\), with compatibility conditions along their boundaries. The precise stratification, metric, and smooth structures required here are not specified in the supplied record. BHV tree space refers to Billera--Holmes--Vogtmann tree space, constructed from Euclidean orthants of nonnegative tree-edge lengths by identifying faces corresponding to edge contractions. The record does not specify the leaf set, rooting convention, or treatment of pendant-edge lengths.

For a curve \(\gamma:I\to X\), where \(I\subseteq\R\) is an interval, write \(\mathcal D_xX\) provisionally for the space of directions to be transported at \(x\). On a smooth stratum, its tangent space is a possible such space; the source does not identify the intended replacement at singular points. A transport construction would have to specify both these direction spaces and maps of the schematic form
\[
\mathcal P^\gamma_{s,t}:\mathcal D_{\gamma(s)}X\longrightarrow\mathcal D_{\gamma(t)}X,
\qquad s,t\in I,
\]
for the parameter pairs to which the construction applies. This notation does not impose linearity, invertibility, preservation of a metric, or any particular composition or compatibility axioms; those requirements are unresolved in the supplied formulation.

\textbf{Problem.} Define parallel transport along sticky principal components in stratified spaces, with BHV tree space as the stated example, so that second principal components and related manifold-statistical techniques can be defined. In particular, specify what directions are transported and how transport is defined when the component lies in a lower-dimensional stratum or passes between strata.

The supplied record does not define a principal component or its stickiness, specify a statistical data set or probability distribution, give an optimization criterion for successive components, or state the axioms that the desired transport must satisfy. Consequently, this is a construction problem whose precise hypotheses and success criteria remain unresolved. The record does not support a universal existence or uniqueness assertion for all stratified spaces, nor does it specify a canonical construction.
\subsection{Short English statement}
The problem asks for a notion of parallel transport along sticky principal components in stratified spaces, especially BHV tree space, that supports second principal components and related statistical constructions. A precise formulation still requires definitions of stickiness, the directions transported at singular points, the transport axioms, and the criterion for successive components. The supplied \texttt{partially\_solved} label is a dated curation assessment, not an independently verified mathematical conclusion.
\subsection{Sources}
\begin{itemize}
\item Local UnsolvedMath registry, supplied complete source record, permanent ID \(30002677\), source problem code \texttt{catalogue entry}, title ``Parallel Transport Along Sticky Principal Components.'' The record belongs to the Oberwolfach Reports Open Problems collection. No individual proposer is supplied; the proposed year is recorded as 2014.
\item Supplied report citation: \emph{Mini-Workshop: Asymptotic Statistics on Stratified Spaces} (2015), \emph{Oberwolfach Reports}, DOI 10.4171/owr/2014/44, pp.\ 2481--2527. \url{https://doi.org/10.4171/owr/2014/44}. The background attributes the problem to workshop open-problem discussions collected in Oberwolfach Reports. No problem-specific page or speaker attribution is supplied.
\item Supplied literature source: ``Pennec, barycentric subspace analysis.'' \url{https://arxiv.org/abs/1607.02833}.
\item Supplied literature source: ``BHV principal components report.'' \url{https://ems.press/content/serial-article-files/46534}.
\item Supplied literature source: ``Oberwolfach report.'' \url{https://doi.org/10.4171/owr/2014/44}.
\item Literature assessment dated exactly \texttt{2026-08-21}: the supplied curation labels the problem partially solved and states that barycentric subspace analysis and multiple-component procedures provide workable stratified-space analogues, while there is no general canonical parallel transport across strata. This is supplied dated curation derived from the OWR material, not independently verified here. The background expressly requires independent verification of time-sensitive status and cited claims.
\item The registry labels the statement \texttt{exact} and describes it as clear, self-contained, publication-ready, and retained byte-for-byte. These are supplied source-verification assertions; the report and cited literature have not been independently consulted for this draft.
\end{itemize}
\end{document}
